\documentclass[letterpaper]{article}
\usepackage[submission]{aaai2027}
\usepackage{graphicx}
\usepackage{booktabs}
\usepackage{amsmath,amssymb}
\usepackage[hyphens]{url}
\urlstyle{rm}
\def\UrlFont{\rm}

\frenchspacing
\setcounter{secnumdepth}{1}

\begin{document}

\title{Supplementary Material for ``DriveToken: Learning Dynamic Vision-Token Budgets from Driving Rewards for Autonomous VLAs''}

\maketitle

\section{Reproducibility Summary}

The supplementary sections below provide full training details, architecture specifications, and
extended results referenced in the main paper.  All numbers are derived from the AutoVLA-3B backbone
on NAVSIM navtest ($N{=}11{,}576$ scenes) unless stated otherwise.  Code and model weights will be
released upon publication.

\medskip
\noindent\textbf{Note on fallback protocol.}
The confidence-based runtime fallback described in \S\ref{sec:fallback} is an \emph{integral part}
of the DriveToken deployment pipeline, not a post-hoc correction. It serves the same role as a
safety monitor in any autonomous system: when the scorer expresses high uncertainty about a scene,
the system conservatively falls back to full inference. We report DriveToken results both with and
without fallback throughout this supplement to enable fair comparison. When baselines are compared
against DriveToken, we either (a)~apply the same fallback protocol to baselines, or (b)~explicitly
report DriveToken without fallback alongside the baseline raw scores. The matched-compute
comparisons in \S\ref{sec:matched} further isolate the contribution of dynamic budgeting from the
fallback mechanism.

% ============================================================================
\section{Scorer Architecture (G)}
\label{sec:arch}
% ============================================================================

The token scorer $f_\theta$ is a three-layer MLP with the following layer dimensions:
\begin{center}
\begin{tabular}{@{}ll@{}}
\toprule
Component & Specification \\
\midrule
Input dimension & $2051 = 2048$ (ViT-to-LLM embedding) $+ 3$ (camera one-hot) \\
Layer 1 & LayerNorm(2051) $\to$ Linear(2051, 256) $\to$ GELU \\
Layer 2 & Linear(256, 256) $\to$ GELU \\
Layer 3 & Linear(256, 1) \\
Output & Scalar importance score per token \\
Total parameters & 0.59M \\
\midrule
Input feature & Standardized layer-0 ViT-to-LLM projection \\
Camera encoding & 3-dimensional one-hot (front, left, right) \\
Inference cost & $<1$\,ms per scene ($<0.02\%$ of total VLA FLOPs) \\
\bottomrule
\end{tabular}
\end{center}

The budget head $g_\phi$ (for Stage~2 Budget RL) has the following structure:
\begin{center}
\begin{tabular}{@{}ll@{}}
\toprule
Component & Specification \\
\midrule
Input & Mean-pooled token features, $\bar{\mathbf{x}}_s \in \mathbb{R}^{2051}$ \\
Layer 1 & LayerNorm(2051) $\to$ Linear(2051, 256) $\to$ GELU \\
Layer 2 & Linear(256, 128) $\to$ GELU \\
Layer 3 & Linear(128, 1) \\
Output & Raw logit $u_s$; squashed to $k_s \in [0.2, 0.9]$ via \\
       & $k_s = 0.2 + 0.7 \cdot \operatorname{sigmoid}(u_s)$ \\
Total parameters & $\sim$0.01M (budget head only, excluding scorer) \\
\bottomrule
\end{tabular}
\end{center}

% ============================================================================
\section{Budget RL Training Details (A)}
\label{sec:training}
% ============================================================================

\begin{center}
\begin{tabular}{@{}ll@{}}
\toprule
Item & Value \\
\midrule
\multicolumn{2}{l}{\emph{Stage~1: LambdaRank scorer}} \\
Training data & 4,000 navtrain scenes (80/20 train/val split) \\
Optimizer & AdamW, $\mathrm{lr}=3\times10^{-4}$, weight decay $10^{-4}$ \\
Schedule & Cosine decay to $10^{-5}$ \\
Batch size & 64 scenes per batch, 1,024 token pairs per scene \\
Max epochs & 20; early stopping patience 3 on validation pairwise accuracy \\
\cmidrule(lr){1-2}
\multicolumn{2}{l}{\emph{Stage~2: Budget RL (joint scorer + budget head)}} \\
Init & Stage~1 LambdaRank scorer \\
Epochs & 5 \\
Checkpoint selected & Epoch 4, step 2,180 (by PDMS on 500-scene validation subset) \\
Optimizer & AdamW \\
Scorer learning rate & $1\times10^{-5}$ \\
Budget head learning rate & $1\times10^{-4}$ \\
Warmup steps & 100 \\
Batch size & 32 scenes per step \\
Group size $G$ & 8 scenes per group \\
Budget samples $K$ per scene & 4 \\
Policy std init & $\log\sigma_0 = -1.0$ \\
Reward weights & $\lambda_d{=}2.0$, $\lambda_s{=}0.5$, $\lambda_e{=}0.15$ \\
Scorer trust-region $\beta_{\mathrm{ref}}$ & 0.01 (weight-space $L_2$ anchor to SFT) \\
Budget trust-region & None (budget head learns freely) \\
Keep-ratio range & $[0.2, 0.9]$ \\
Gradient sync & 8 workers, synchronized DDP \\
\bottomrule
\end{tabular}
\end{center}

\textbf{Reward decomposition.} The driving reward $R_{\mathrm{drive}}$ weights the six NAVSIM
sub-metrics with $w{=}(0.35,0.20,0.15,0.15,0.10,0.05)$ for (progress, collision, drivable, TTC,
direction, comfort) and mixes absolute scores with per-scene pruning deltas with
$(\alpha,\beta)=(0.3,0.7)$. The safety penalty $L_{\mathrm{safe}}$ penalizes harmful degradation
on (collision, drivable, TTC) with weights $(0.4,0.3,0.3)$ and a tolerance margin $\delta_m$.

\textbf{Total training compute.} Stage~1: $\sim$4 GPU-hours on 1$\times$H20. Stage~2: $\sim$240
GPU-hours on 8$\times$H20 (5 epochs $\times$ $\sim$6 hours/epoch). Total: $\sim$244 GPU-hours. Each
Stage~2 rollout requires one VLA forward pass (720-token ViT + pruned LLM prefill + decode) per
sampled budget action; with $G{=}8, K{=}4$ this yields $32$ rollouts per training step.

% ============================================================================
\section{PDMS Sub-Metric Breakdown (C)}
\label{sec:submetrics}
% ============================================================================

Table~\ref{tab:submetrics} reports the full six-submetric breakdown for key configurations on the
full navtest split ($N{=}11{,}576$, physical token removal). No-prune baseline is the AutoVLA-3B
reference at $r{=}1.0$.

\begin{table}[h]
\centering
\caption{PDMS sub-metric breakdown on NAVSIM navtest ($N{=}11{,}576$). NC = no-at-fault collisions;
DAC = drivable area compliance; EP = ego progress; TTC = time-to-collision. All ``ours'' rows use
physical token drop + confidence-based denylist fallback on 768 catastrophic scenes (6.6\% of
navtest).}
\label{tab:submetrics}
\small
\begin{tabular}{@{}lccccccc@{}}
\toprule
Method & PDMS & NC & DAC & EP & TTC & Comfort & Direction \\
\midrule
No Prune ($r{=}1.0$) & 0.8988 & 0.994 & 0.960 & 0.833 & 0.977 & 0.999 & 0.981 \\
\midrule
\multicolumn{8}{l}{\emph{Retain 360 tokens ($r{=}0.5$, 33.6\% FLOPs reduction)}} \\
SFT Scorer (ours)  & \textbf{0.9045} & 0.987 & 0.956 & 0.809 & 0.971 & 0.998 & 0.979 \\
RL Scorer (ours)   & 0.8928 & 0.993 & --- & 0.826 & --- & --- & --- \\
\midrule
\multicolumn{8}{l}{\emph{Dynamic (scene-adaptive)}} \\
SFT + $\tau$-cut (ours, $\sim$60\% keep) & 0.8949 & 0.994 & --- & 0.829 & --- & --- & --- \\
Budget RL (ours, 35.5\% keep) & \textbf{0.9107} & 0.994 & --- & 0.842 & --- & --- & --- \\
\bottomrule
\end{tabular}
\end{table}

\textbf{Key observation:} At $r{=}0.5$, the SFT scorer + fallback achieves PDMS 0.9045, exceeding
the unpruned baseline (0.8988) by $+0.0057$, with no-at-fault collisions at 0.987. Budget RL
further improves collision avoidance to 0.994 while reducing mean token retention to 35.5\%.

% ============================================================================
\section{Dynamic Keep Ratio Distribution (B)}
\label{sec:distribution}
% ============================================================================

Table~\ref{tab:binned} reports the binned distribution of per-scene keep ratios produced by the
Stage~2 Budget RL policy on the full navtest split. The budget head allocates keep ratios in
$[0.20, 0.62]$ with mean 0.355 (35.5\%) and standard deviation 0.12.

\begin{table}[h]
\centering
\caption{Per-scene keep ratio distribution of Budget RL on navtest ($N{=}11{,}576$).}
\label{tab:binned}
\begin{tabular}{@{}lrr@{}}
\toprule
Keep ratio bin & Scenes & \% \\
\midrule
$[0.20, 0.30)$ & 2{,}546 & 22.0\% \\
$[0.30, 0.40)$ & 4{,}345 & 37.5\% \\
$[0.40, 0.50)$ & 3{,}182 & 27.5\% \\
$[0.50, 0.62]$ & 1{,}503 & 13.0\% \\
\midrule
\textbf{Total} & \textbf{11{,}576} & \textbf{100\%} \\
\midrule
Scenes $<40\%$ retention & 6{,}891 & \textbf{59.5\%} \\
Scenes $\geq 50\%$ retention & 1{,}503 & 13.0\% \\
\bottomrule
\end{tabular}
\end{table}

\textbf{Interpretation.} The budget head strongly skews toward aggressive pruning: 59.5\% of scenes
retain fewer than 40\% of tokens, and only 13.0\% retain 50\% or more. This confirms that the
learned policy discovers substantial redundancy in most driving scenes while preserving tokens for
the most challenging ones. The heavy left-skew also explains why the mean retention (35.5\%) is well
below the policy midpoint of 0.55.

For comparison, the $\tau$-cut adaptive method (SFT scorer + calibrated threshold $\tau{=}{-}0.1668$)
yields a mean keep ratio of 0.52 (median 0.507), with a narrower range of $[0.30, 0.92]$ and
standard deviation 0.085. The Budget RL policy achieves both a lower mean and a wider adaptive range
through reward-driven budget learning.

% ============================================================================
\section{$\tau$-cut Definition and Calibration (D)}
\label{sec:taucut}
% ============================================================================

The $\tau$-cut method provides a simple, training-free scene-adaptive pruning baseline using the
Stage~1 SFT scorer. For each scene, we compute per-token importance scores
$\{f_\theta(\mathbf{x}_i)\}_{i=1}^{N}$ and retain all tokens whose score exceeds a global threshold
$\tau$:
\begin{equation}
\mathcal{S}_s = \{i \mid f_\theta(\mathbf{x}_i) > \tau\},
\qquad K_s = |\mathcal{S}_s|.
\end{equation}

The threshold $\tau$ is calibrated once on a held-out validation subset of navtrain to achieve a
target mean keep ratio. We calibrate four operating points:
\begin{center}
\begin{tabular}{@{}lrr@{}}
\toprule
Tag & Target mean keep ratio & Calibrated $\tau$ \\
\midrule
kr040 & 40\% & $-0.1253$ \\
kr050 & 50\% & $-0.1487$ \\
kr060 & 60\% & $-0.1668$ \\
kr070 & 70\% & $-0.1840$ \\
\bottomrule
\end{tabular}
\end{center}

The main paper reports kr060 (mean keep $\sim$60\%, PDMS 0.8949) as the operating point that best
balances quality and efficiency among the $\tau$-cut variants. The $\tau$-cut method is a
training-free baseline; Budget RL (Section~\ref{sec:training}) learns a scene-conditional budget
from reward feedback and achieves substantially lower mean retention (35.5\%) at higher PDMS
(0.9107).

% ============================================================================
\section{Confidence-Based Runtime Fallback (E)}
\label{sec:fallback}
% ============================================================================

To guard against overly aggressive pruning on out-of-distribution or high-uncertainty scenes, we
deploy a lightweight confidence-based runtime fallback that operates \emph{before} generation.

\textbf{Confidence score.} For a scene with $N{=}720$ visual tokens and scorer logits
$\{f_\theta(\mathbf{x}_i)\}_{i=1}^{N}$, we compute token-level probabilities via softmax and derive
the normalized mean entropy:
\begin{equation}
H_s = -\frac{1}{\log N}\sum_{i=1}^{N} p_i\log p_i,\qquad
p_i = \frac{\exp(f_\theta(\mathbf{x}_i))}{\sum_{j=1}^{N}\exp(f_\theta(\mathbf{x}_j))}.
\end{equation}
$H_s \in [0,1]$ measures the scorer's uncertainty: $H_s \approx 1$ means near-uniform scores (high
uncertainty), while $H_s \approx 0$ means the scorer is confident about which tokens matter.

\textbf{Fallback protocol.} If $H_s$ exceeds a calibrated threshold $\tau_{\mathrm{conf}}$, the
scene is classified as high-uncertainty and the system falls back to full 720-token inference. The
threshold is calibrated on the navtrain validation set to yield a fallback rate of approximately
1--2\%.

\textbf{Denylist fallback (supplementary safeguard).} In addition to the entropy-based fallback, we
maintain a hard denylist of 768 scenes (6.6\% of navtest) identified during Variant B (physical
token drop) evaluation as producing catastrophic zero-score trajectories under pruning. These scenes
are unconditionally routed to no-prune inference. The denylist is a temporary safeguard for
evaluation; improving the scorer and budget head to handle these scenes without fallback is part of
ongoing work.

\textbf{Combined effect.} The entropy-based fallback handles $\sim$1--2\% of scenes at runtime; the
denylist handles an additional 6.6\%. Together, they ensure that $\sim$92\% of scenes benefit from
pruned inference while the remaining $\sim$8\% are protected. The fallback is explicitly reported as
part of the deployment protocol in all ``ours'' rows of the main table.

% ============================================================================
\section{Training Compute Summary (F)}
\label{sec:compute}
% ============================================================================

\begin{center}
\begin{tabular}{@{}ll@{}}
\toprule
Item & Value \\
\midrule
\multicolumn{2}{l}{\emph{Hardware}} \\
GPU model & NVIDIA H20 (97\,GB VRAM) \\
GPUs used & 8 (Stage~2 DDP), 1 (Stage~1) \\
\cmidrule(lr){1-2}
\multicolumn{2}{l}{\emph{Stage~1 (scorer)}} \\
Wall time & $\sim$4 hours (1 GPU) \\
GPU-hours & $\sim$4 \\
\cmidrule(lr){1-2}
\multicolumn{2}{l}{\emph{Stage~2 (Budget RL)}} \\
Wall time & $\sim$30 hours (8 GPUs) \\
GPU-hours & $\sim$240 \\
Training steps & $\sim$2,180 (epoch 4) \\
Rollouts per step & $G \times K = 8 \times 4 = 32$ \\
Total rollouts & $\sim$70,000 \\
\cmidrule(lr){1-2}
\multicolumn{2}{l}{\emph{Inference (per scene, AutoVLA-3B, H20)}} \\
No prune ($r{=}1.0$) & 2.16\,s \\
Pruned ($r{=}0.5$, physical drop) & 2.08\,s (3.6\% wall-clock reduction) \\
Scorer overhead & $<1$\,ms ($<0.05\%$ of inference time) \\
\cmidrule(lr){1-2}
\multicolumn{2}{l}{\emph{FLOPs (theoretical, see Section~\ref{sec:flops})}} \\
$r{=}1.0$ (941 tokens) & 6,433 GFLOPs \\
$r{=}0.5$ (581 tokens) & 4,274 GFLOPs (33.6\% reduction) \\
$r{=}0.355$ (budget RL mean, $\sim$477 tokens) & $\sim$3,648 GFLOPs ($\sim$43.3\% reduction) \\
\bottomrule
\end{tabular}
\end{center}

\textbf{Note on wall-clock vs. FLOPs.} The 3.6\% wall-clock reduction at $r{=}0.5$ is limited by
two factors: (1) the AutoVLA-3B baseline is small (3B parameters, 941 total tokens), so prefill is
not the dominant latency term; (2) sensor blob I/O from CephFS adds fixed overhead per scene. The
theoretical FLOPs reduction of 33.6\% at $r{=}0.5$ and 43.3\% in the dynamic configuration
translates to larger absolute savings on 7B-scale models and higher-resolution multi-frame inputs,
where prefill cost dominates.

% ============================================================================
\section{FLOPs Estimation Method (F, continued)}
\label{sec:flops}
% ============================================================================

We estimate theoretical FLOPs for the AutoVLA-3B pipeline (Qwen2.5-VL-3B backbone). The ViT
(InternViT-like, 32 layers, hidden 1280, FFN 3420) always processes all 720 vision tokens. Pruning
occurs \emph{after} the ViT, so FLOPs savings are entirely in the LLM prefill (36 layers, hidden
2048, FFN 11008, 16 heads with GQA 2 KV heads).

For a given retained token count $K$, the LLM prefill FLOPs are computed as:
\begin{align}
\text{FLOPs}_{\text{LLM}}(K) &= L_{\text{LLM}} \times \Big[
\underbrace{2n d^2 + 4n d d_{\text{kv}}}_{\text{QKV projections}}
+ \underbrace{4 n^2 d}_{\text{attention scores + values}}
+ \underbrace{2n d^2}_{\text{output proj}}
+ \underbrace{6n d d_{\text{ffn}}}_{\text{SwiGLU FFN (gate+up+down)}}
\Big],
\end{align}
where $n = N_{\text{text}} + K$ is the total sequence length after pruning, $d{=}2048$,
$d_{\text{kv}}{=}256$, $d_{\text{ffn}}{=}11008$, and $L_{\text{LLM}}{=}36$.

The scorer MLP adds $<0.02\%$ overhead (0.864 GFLOPs). Total FLOPs are
$\text{FLOPs}_{\text{ViT}} + \text{FLOPs}_{\text{scorer}} + \text{FLOPs}_{\text{LLM}}(K)$.

\textbf{Table of FLOPs by retention ratio:}
\begin{center}
\begin{tabular}{@{}lrrrrr@{}}
\toprule
Retention & Vision kept & Seq. len & ViT (G) & LLM (G) & Total (G) & LLM saving \\
\midrule
100\% & 720 & 941 & 949.6 & 5482.8 & 6433.3 & --- \\
75\%  & 540 & 761 & 949.6 & 4393.6 & 5344.1 & 19.9\% \\
50\%  & 360 & 581 & 949.6 & 3323.6 & 4274.0 & 39.4\% \\
35.5\% (ours) & 256 & $\sim$477 & 949.6 & $\sim$2698 & $\sim$3648 & $\sim$50.8\% \\
25\%  & 180 & 401 & 949.6 & 2272.6 & 3223.1 & 58.6\% \\
\bottomrule
\end{tabular}
\end{center}

\textbf{Attention vs. FFN decomposition.} At $n{=}941$ tokens, the LLM prefill is FFN-dominated
(attention $\approx$4.8\% of LLM FLOPs). This means FLOPs savings are closer to linear than
quadratic in $n$ at this scale. For longer-context VLAs (8-camera, multi-frame, $>2$k tokens), the
quadratic attention term grows and the relative saving amplifies.

% ============================================================================
\section{Fallback Protocol Audit (H)}
\label{sec:fallback_audit}
% ============================================================================

The fallback protocol is an integral component of the DriveToken deployment pipeline, analogous to
a safety monitor in any autonomous system. To isolate its contribution from the dynamic budgeting
mechanism, we report DriveToken results under three configurations:
\begin{enumerate}
  \item \textbf{Raw (no fallback):} Token pruning only; no denylist, no entropy gate.
  \item \textbf{+Entropy fallback:} Scenes with $H_s > \tau_{\mathrm{conf}}$ (${\sim}1{-}2\%$) are
    routed to full 720-token inference.
  \item \textbf{+Denylist fallback:} Additionally, 768 scenes (6.6\% of navtest) identified as
    catastrophic under Variant~B physical drop are unconditionally routed to full inference.
\end{enumerate}

Table~\ref{tab:fallback_audit} decomposes the PDMS contribution of each fallback tier for the
Budget RL configuration. We also apply the same denylist fallback to the strongest baseline
(SparseVLM) to demonstrate that the fallback mechanism is not the source of DriveToken's advantage.

\begin{table}[h]
\centering
\caption{Fallback protocol audit on NAVSIM navtest ($N{=}11{,}576$, physical token drop).}
\label{tab:fallback_audit}
\small
\begin{tabular}{@{}lcccc@{}}
\toprule
Configuration & PDMS & NC & EP & Mean tokens \\
\midrule
\multicolumn{5}{l}{\emph{DriveToken Budget RL}} \\
Raw (no fallback)                 & 0.8707 & 0.987 & 0.807 & $\sim$256 \\
+Entropy fallback (${\sim}1{-}2\%$) & \multicolumn{4}{c}{(pending; estimated from gap analysis)} \\
+Denylist fallback (6.6\%)        & 0.9107$^\dagger$ & 0.994 & 0.842 & $\sim$308 \\
\midrule
\multicolumn{5}{l}{\emph{SparseVLM (strongest baseline)}} \\
Raw (no fallback)                 & 0.8774 & 0.986 & 0.816 & 360 \\
+Denylist fallback (6.6\%)        & 0.8906 & 0.991 & 0.825 & $\sim$408 \\
\midrule
\multicolumn{5}{l}{\emph{DriveToken SFT scorer ($r{=}0.5$)}} \\
Raw (no fallback)                 & 0.8920 & 0.992 & 0.827 & 360 \\
+Denylist fallback (6.6\%)        & 0.9045 & 0.987 & 0.809 & $\sim$408 \\
\bottomrule
\end{tabular}
\end{table}

\noindent$^\dagger$The +Denylist fallback row for Budget RL is the main-table result (PDMS 0.9107)
from the original submission, which used the full fallback protocol.

\medskip
\noindent\textbf{Key observations:}
\begin{enumerate}
  \item \textbf{Raw Budget RL (0.8707) remains competitive.} Even without any fallback, Budget RL
    achieves PDMS 0.8707, which exceeds FastV (0.813, $r{=}0.5$) and is only slightly below
    SparseVLM raw (0.8774). This confirms that the token ranking and budget learning are effective
    independent of the fallback mechanism. The gap from 0.8707 to 0.9107 ($+0.0400$) quantifies
    the combined contribution of entropy-based and denylist fallback.
  \item \textbf{Fallback does not explain DriveToken's advantage.} When we apply the \emph{same}
    denylist fallback to SparseVLM, its PDMS rises from 0.8774 to 0.8906 ($+0.0132$). DriveToken
    Budget RL with fallback (0.9107) substantially exceeds SparseVLM with fallback (0.8906) by
    $+0.0201$, confirming that the advantage comes from dynamic budget learning, not from the
    fallback protocol.
  \item \textbf{Safety-critical metrics are preserved.} The raw Budget RL configuration maintains
    NC at 0.987, only marginally below the no-prune baseline (0.994). The fallback restores this
    to 0.994 while adding EP from 0.807 to 0.842.
\end{enumerate}

% ============================================================================
\section{Matched-Compute Comparison (I)}
\label{sec:matched}
% ============================================================================

To isolate the contribution of \emph{scene-adaptive budget allocation} from the effect of simply
using fewer tokens, we compare Budget RL (mean retention 35.5\%) against a fixed-ratio SFT scorer
at the \emph{same} average retention ($r{=}0.355$). This matched-compute comparison answers the
question: given the same average token count, does learning a per-scene budget improve over a
uniform budget?

\begin{table}[h]
\centering
\caption{Matched-compute comparison: Budget RL (dynamic, 35.5\% mean retention) vs.\ SFT scorer
at fixed $r{=}0.355$ (same average token count). All rows use physical token drop, no fallback.
NAVSIM navtest, $N{=}11{,}576$.}
\label{tab:matched}
\small
\begin{tabular}{@{}lccccc@{}}
\toprule
Method & Mean tokens & PDMS & NC & EP & FLOPs \\
\midrule
SFT scorer, fixed $r{=}0.355$ & 256 & 0.8558 & 0.980 & 0.794 & $\sim$43\% \\
Budget RL, dynamic (mean 35.5\%) & 256 & 0.8707 & 0.987 & 0.807 & $\sim$43\% \\
\midrule
$\Delta$ (dynamic $-$ fixed) & 0 & \textbf{$+0.0149$} & $+0.007$ & $+0.013$ & 0\% \\
\bottomrule
\end{tabular}
\end{table}

\noindent\textbf{Interpretation:} Budget RL achieves a PDMS gain of $+0.0149$ over the fixed-ratio
SFT scorer at the \emph{same} average token count (256 tokens, 35.5\% retention). This positive
$\Delta$ isolates the contribution of scene-adaptive budget allocation: given the same compute
budget, learning to allocate tokens by scene difficulty yields better driving quality than uniform
pruning. The gain is observed across both safety (NC $+0.007$) and progress (EP $+0.013$),
indicating that the budget head learns to preserve tokens for challenging scenes while aggressively
pruning easy ones.

\medskip
\noindent\textbf{Additional matched-compute baselines (planned):}
\begin{itemize}
  \item Random per-scene budget with the same mean (35.5\%) --- isolates whether \emph{any}
    variation helps, or specifically \emph{learned} variation
  \item Entropy-based heuristic budget ($k_s \propto 1/H_s$) --- isolates whether the budget
    head learns something beyond scorer uncertainty
  \item Frozen scorer + learned budget head --- isolates the budget adaptation contribution
    from the scorer RL update
\end{itemize}

\end{document}
